\documentclass[11pt,letterpaper]{article}
\usepackage{physical_ai_legislation}
\usepackage[backend=biber,style=numeric,sorting=none]{biblatex}
\addbibresource{physical_ai_trial_authorization.bib}

\title{%
  \textbf{Senate Bill 1042}\\[6pt]
  \large California Physical AI Oncology Clinical Trial\\
  Authorization and Site Establishment Act of 2026\\[6pt]
  \normalsize An act to add Chapter 4.7 (commencing with Section 1399.900)\\
  to Division 2 of the Health and Safety Code,\\
  relating to physical artificial intelligence in oncology clinical trials.
}
\author{%
  Kevin Kawchak\\
  CEO, ChemicalQDevice\\
  \texttt{kevink@chemicalqdevice.com}
}
\date{March 24, 2026}

\begin{document}
\maketitle
\thispagestyle{fancy}

\begin{center}
\footnotesize
\textit{Unofficial independent draft derived from prior legislation;
no third-party endorsement, sponsorship, affiliation, or authorization
is expressed or implied; and was adapted using Claude Code Opus 4.6.}
\end{center}

\vspace{0.5em}

\begin{center}
\footnotesize
\textit{This work was adapted from original CFR documents in the public domain.
The original ICH document is copyrighted and may be used, reproduced,
incorporated into other works, adapted, modified, translated or distributed
under a public license. This current work is not endorsed or sponsored by
CFR, ICH, or FDA; and was adapted using Claude Code Opus 4.6.}
\end{center}

\vspace{1em}
\tableofcontents
\newpage

\section{Legislative Findings and Declarations}

The Legislature finds and declares all of the following:

\begin{enumerate}[label=(\alph*)]
\item Physical artificial intelligence, including robotic surgical systems,
  collaborative robots, radiotherapy positioning systems, imaging assistants,
  rehabilitation exoskeletons, and companion robots, represents a transformative
  advancement in oncology clinical care delivery.

\item A 24-hour on-demand Physical AI oncology clinical trial simulation
  conducted on March 23, 2026, demonstrated that a single facility equipped
  with 10 robot types and 29 robot instances successfully served 168 unique
  patients across 15 cancer types in 24 hours, compared to 5 to 15 patients per
  day in traditional trial formats~\cite{newtrial2026}.

\item The simulation achieved a cumulative site Physical AI Standard Level (PSL)
  score of 63.4 to 64.4, within the Advanced Site band, with 99.7 percent
  facility uptime and an average patient wait time of 8 minutes.

\item All seven adverse events documented during the simulation were Grade 1 or
  Grade 2, all were resolved within the same hour they occurred, and zero
  patient harm resulted, demonstrating that Physical AI systems can manage
  both minute-level clinical details and sustained 24-hour
  operations~\cite{newtrial2026}.

\item The Physical AI trial format shifts control toward the patient by enabling
  patient-chosen scheduling on a 24/7 basis, increasing patient throughput by a
  factor of 10, reducing per-patient costs by 75 to 85 percent, and improving
  retention rates from 70 to 80 percent in traditional formats to 85 to
  95 percent~\cite{formatcomparison2026}.

\item Staffing requirements are reduced from 80 to 120 full-time equivalents in
  traditional trial sites to 8 to 10 full-time equivalents in a Physical AI
  trial site, with the majority of clinical procedures, monitoring,
  documentation, and regulatory compliance tasks performed
  autonomously~\cite{sitespec2026}.

\item The Unification Standard Level (USL) framework provides a standardized
  1.0 to 10.0 scoring system for evaluating physical AI robot readiness across
  four dimensions: simulation framework switching, AI integration, cross-robot
  progress sharing, and clinical trial collaboration~\cite{usl2026}.

\item Three adapted regulatory frameworks, covering ICH E6(R3) good clinical
  practice, 21 CFR Part 50 human subjects protection, and 21 CFR Part 312
  investigational new drug applications, provide comprehensive governance for
  Physical AI oncology trials~\cite{iche6r3adapt,cfr50adapt,cfr312adapt}.

\item The cancer patient journey through a regulated Physical AI oncology trial
  has been documented across 10 sequential stages from pre-screening through
  data lock, demonstrating complete Response outcomes and a reduction in
  recurrence risk from 35 percent to 3 percent~\cite{patientjourney2026}.

\item California has an opportunity to lead the nation in establishing safe,
  regulated, and patient-centric Physical AI oncology clinical trial sites that
  serve more patients with more robots and fewer workers while maintaining the
  highest standards of clinical care.
\end{enumerate}

\section{Definitions}

For purposes of this chapter, the following definitions apply:

\begin{enumerate}[label=(\alph*)]
\item ``Physical AI system'' means a robotic or embodied artificial intelligence
  system that physically interacts with patients, clinical specimens, or medical
  equipment in the context of an oncology clinical trial. Physical AI systems
  include, but are not limited to, the following 10 categories:
  \begin{enumerate}[label=(\arabic*)]
  \item Surgical robots that perform tumor resection through minimally invasive
    procedures.
  \item Collaborative robots (cobots) that stabilize patient limbs and assist
    with needle biopsy procedures.
  \item Radiotherapy patient-positioning robots that align patients with
    submillimeter precision for radiation delivery.
  \item Robotic needle-placement systems that use imaging guidance to place
    biopsy or ablation needles with high accuracy.
  \item Social companion robots that provide emotional support and
    anxiety reduction for patients, including pediatric patients.
  \item Humanoid robots that conduct rehabilitation exercises and patient
    engagement activities.
  \item Radiotherapy motion-management and tracking robots that monitor
    patient breathing and gate radiation beams accordingly.
  \item Imaging assistant robots that perform automated ultrasound and
    diagnostic imaging procedures.
  \item Steerable needle and needle-steering robots that guide flexible
    needles through tissue for targeted ablation.
  \item Rehabilitation exoskeletons and robotic gait trainers that support
    post-surgical mobility recovery.
  \end{enumerate}

\item ``Physical AI Standard Level'' or ``PSL'' means the scoring framework
  that evaluates each robot type on three regulatory dimensions: Omniscient
  (ICH E6(R3) compliance), Omnipresent (21 CFR Part 50 compliance), and
  Omnipotent (21 CFR Part 312 compliance), with each dimension scored from
  0.0 to 10.0 per robot type.

\item ``Unification Standard Level'' or ``USL'' means the 1.0 to 10.0
  scoring framework evaluating physical AI robot readiness across four equally
  weighted dimensions: simulation framework switching, AI integration,
  cross-robot progress sharing, and clinical trial
  collaboration~\cite{usl2026}.

\item ``On-demand trial format'' means a clinical trial structure in which
  patients may schedule procedures at any time during a 24-hour period, 7 days
  per week, with Physical AI systems providing continuous clinical care.

\item ``Digital twin'' means a patient-specific computational model calibrated
  to individual anatomy and treatment parameters, used for procedure planning,
  rehearsal, and intraoperative guidance.

\item ``MCP-PAI standard'' means the Model Context Protocol for Physical AI,
  a five-server interoperability topology providing authorization, FHIR, DICOM,
  ledger, and provenance services with hash-chained audit trails.

\item ``Qualified site'' means a facility that has met all building, premises,
  parking, staffing, equipment, and regulatory requirements established under
  this chapter and has received authorization from the department to conduct
  Physical AI oncology clinical trials.

\item ``Department'' means the California Department of Public Health.

\item ``Sim-to-real validation'' means the process of verifying that a Physical
  AI system's performance in simulated environments transfers to real clinical
  settings within specified tolerances, including trajectory deviation below
  2 millimeters and force discrepancy below 0.5 Newtons.
\end{enumerate}

\section{Authorization to Establish Physical AI\\Oncology Clinical Trial Sites}

\subsection{Site Authorization}

\begin{enumerate}[label=(\alph*)]
\item The department shall establish a process for authorizing Physical AI
  oncology clinical trial sites in California, beginning with the first site
  in a prominent and safe San Francisco location.

\item A qualified entity may apply to the department for authorization to
  establish and operate a Physical AI oncology clinical trial site in a new
  building with dedicated parking facilities.

\item The authorization process shall include review of the following:
  \begin{enumerate}[label=(\arabic*)]
  \item Facility design and construction plans meeting the Physical AI
    oncology trial facility building code standards established under
    this chapter.
  \item Equipment specifications for all 10 Physical AI system categories
    to be deployed at the site.
  \item Staffing plans demonstrating adequate human oversight with a minimum
    of 8 to 10 full-time equivalent positions across three shifts.
  \item USL scores for each Physical AI system, with minimum thresholds of
    7.0 for surgical systems, 6.0 for therapeutic positioning systems, 6.0 for
    diagnostic systems, 4.0 for rehabilitative systems, and 3.0 for companion
    systems.
  \item Phase 0 simulation validation results demonstrating cross-framework
    consistency using two or more simulation frameworks.
  \item Compliance plans for ICH E6(R3), 21 CFR Part 50, and 21 CFR Part 312
    as adapted for Physical AI systems.
  \item Emergency response and cybersecurity plans.
  \item Parking and patient transportation accessibility plans.
  \end{enumerate}

\item The department shall act on a complete application within 120 days of
  receipt.
\end{enumerate}

\subsection{Statewide Applicability}

\begin{enumerate}[label=(\alph*)]
\item The authorization requirements and standards established under this
  chapter shall apply uniformly to Physical AI oncology clinical trial sites
  throughout the State of California.

\item A city, county, or city and county may adopt additional requirements for
  Physical AI oncology clinical trial sites, provided those requirements do not
  conflict with or impose lesser standards than those established under this
  chapter.

\item The department shall publish model guidelines to assist local
  jurisdictions in integrating Physical AI oncology clinical trial site
  requirements into existing land use, building, and health facility
  permitting processes.
\end{enumerate}

\section{Physical AI System Requirements}

\subsection{USL and PSL Compliance}

\begin{enumerate}[label=(\alph*)]
\item Each Physical AI system deployed at an authorized site shall maintain a
  USL score at or above the minimum threshold for its system category as
  specified in Section 3.1(c)(4).

\item The site shall maintain a cumulative PSL score that places it within the
  Advanced Site band (60.0 or above on the 0 to 100 scale) as a condition of
  continued authorization.

\item PSL scores shall be calculated and reported monthly using the three
  regulatory dimensions established in the PSL framework, with changes limited
  to a maximum of 0.3 per hour per dimension per robot type to ensure scoring
  stability~\cite{pslframework2026}.

\item A site whose cumulative PSL score falls below the Advanced Site band for
  two consecutive monthly reporting periods shall submit a corrective action
  plan to the department within 30 days.
\end{enumerate}

\subsection{Simulation Validation Requirements}

\begin{enumerate}[label=(\alph*)]
\item Before deployment at a qualified site, each Physical AI system shall
  complete Phase 0 simulation validation using two or more validated
  frameworks, which may include NVIDIA Isaac Lab, MuJoCo, Gazebo Sim, and
  PyBullet~\cite{iche6r3adapt}.

\item Cross-framework validation shall demonstrate trajectory deviation below
  2 millimeters and force discrepancy below 0.5 Newtons.

\item Digital twin models shall achieve a minimum fidelity of 90 percent when
  compared to patient-specific imaging and anatomical data.

\item Failure mode analysis shall be conducted across a minimum of 1,000
  simulated procedures per Physical AI system type before clinical
  deployment~\cite{cfr312adapt}.
\end{enumerate}

\subsection{Informed Consent and Patient Rights}

\begin{enumerate}[label=(\alph*)]
\item Physical AI oncology clinical trial sites shall comply with all informed
  consent requirements established in the adapted 21 CFR Part 50 for Physical
  AI systems, including disclosure of robotic system roles, USL scores in plain
  language, safety records, and emergency stop provisions~\cite{cfr50adapt}.

\item Patients shall have the right to schedule procedures at any time during
  the 24-hour operating period, consistent with the on-demand trial format.

\item Patients shall have the right to request human intervention or oversight
  at any point during a Physical AI-assisted procedure.

\item Pediatric patients shall receive age-appropriate robotic demonstrations
  and assent procedures as specified in the adapted 21 CFR Part 50
  \S~50.50 through \S~50.56~\cite{cfr50adapt}.

\item Each patient shall be informed of the specific robots that will be
  involved in their care, the role of each robot (autonomous, assistive, or
  shared-autonomy), and the qualifications of human oversight
  personnel~\cite{patientinstructions2026}.
\end{enumerate}

\section{Safety and Reporting Requirements}

\subsection{Adverse Event Reporting}

\begin{enumerate}[label=(\alph*)]
\item The site shall report all Physical AI-related adverse events to the
  department and to the Food and Drug Administration in accordance with the
  timelines established in the adapted 21 CFR Part 312 \S~312.32, including
  7-day reporting for fatal or life-threatening events and 15-day reporting for
  serious events~\cite{cfr312adapt}.

\item Seven categories of reportable Physical AI adverse events shall be
  tracked: serious Physical AI adverse events, Physical AI-drug interaction
  events, cybersecurity incidents, AI model performance degradation,
  sim-to-real divergence, digital twin discrepancy, and safety review committee
  findings.

\item The site shall maintain a Physical AI Safety Review Committee with
  expertise in robotic engineering, AI and machine learning, oncology, and
  cybersecurity, meeting at intervals not exceeding 90 days.
\end{enumerate}

\subsection{Data Protection and Audit Trails}

\begin{enumerate}[label=(\alph*)]
\item All Physical AI systems shall maintain hash-chained audit trails using
  SHA-256 or equivalent cryptographic standards, as required by the adapted
  ICH E6(R3) and 21 CFR Part 11~\cite{iche6r3adapt}.

\item Patient data processed through MCP-PAI servers shall be de-identified
  using HIPAA Safe Harbor methodology with deny-by-default authorization
  policies.

\item The site shall implement federated learning protocols with differential
  privacy protections for multi-site data sharing, as validated in the
  simulation where 70 federated learning rounds across 12 sites used secure
  multi-party computation with differential privacy at epsilon equals
  1.0~\cite{patientjourney2026}.
\end{enumerate}

\section{Facility and Infrastructure Standards}

\begin{enumerate}[label=(\alph*)]
\item Physical AI oncology clinical trial sites shall be constructed in
  accordance with the building code standards established under this chapter,
  including a minimum floor area of 85,000 to 100,000 square feet, HEPA
  filtration in surgical suites, shielded radiotherapy vaults, and dedicated
  server rooms for AI computation infrastructure~\cite{sitespec2026}.

\item The site shall provide dedicated parking facilities meeting the standards
  established under this chapter, including accessible parking, patient
  drop-off zones, and emergency vehicle access.

\item Power infrastructure shall support 800 to 1,200 kilowatts of sustained
  electrical load with redundant power supply and uninterruptible power systems
  for all critical clinical areas.

\item Network infrastructure shall provide a minimum of 10 gigabits per second
  internal bandwidth with cybersecurity protections meeting the requirements of
  the MCP-PAI standard~\cite{sitespec2026}.
\end{enumerate}

\section{Workforce Transition and Training}

\begin{enumerate}[label=(\alph*)]
\item The department, in consultation with the Employment Development Department
  and the California Workforce Development Board, shall establish programs to
  assist clinical trial workers in transitioning to roles in Physical AI trial
  sites or related fields.

\item Each authorized site shall provide comprehensive training for all human
  staff in Physical AI system operation, emergency response, patient
  interaction, and regulatory compliance.

\item Training requirements shall include demonstrated competency in
  understanding Physical AI system capabilities, limitations, and failure modes,
  as well as the ability to intervene manually when necessary.

\item The Legislature recognizes that while Physical AI trial sites require
  fewer workers (8 to 10 full-time equivalents compared to 80 to 120 in
  traditional formats), the transition enables more patients to be treated,
  with the simulation demonstrating 168 patients served in 24 hours
  compared to approximately 5 to 15 per day traditionally.
\end{enumerate}

\section{Relationship to Existing California Law}

This chapter is intended to complement and operate in conjunction with
existing California laws governing health care and artificial intelligence,
including but not limited to the following:

\begin{enumerate}[label=(\alph*)]
\item AB 489 (2025), concerning health care professions and deceptive terms
  related to artificial intelligence, which ensures clear communication about
  AI involvement in clinical care.

\item AB 3030 (2024), concerning health care services and artificial
  intelligence, which establishes transparency requirements for AI in health
  care delivery.

\item SB 1120 (2024), concerning health care coverage and utilization review,
  which addresses AI-assisted clinical decision-making.

\item SB 243 (2025), concerning companion chatbots, which establishes safety
  standards for AI companion systems relevant to social companion robots used
  in Physical AI oncology trials.

\item AB 2013 (2024), concerning generative artificial intelligence and
  training data transparency, which applies to the generative AI and VLA
  models integrated into Physical AI systems.
\end{enumerate}

\section{Appropriation}

The sum of fifty million dollars (\$50,000,000) is hereby appropriated from the
General Fund to the department for the purposes of implementing this chapter,
including the establishment of the site authorization process, development of
model guidelines, workforce transition programs, and ongoing oversight of
authorized Physical AI oncology clinical trial sites.

\section{Effective Date}

This act shall take effect on January 1, 2027.

\printbibliography[title={References}]

\end{document}
